\subsection{Comparison with the rule-based baseline}

To assess the practical contribution of LLM-assisted extraction, the same
human-annotated reference sample was also evaluated with a deterministic
rule-based baseline. The comparison uses the same 60 documents and the same
gold labels as the LLM evaluation for \texttt{public\_vise},
\texttt{echelle}, and \texttt{nature\_initiative}. It is therefore a paired
comparison: differences cannot be attributed to a different sample or to a
different reference annotation.

For \texttt{public\_vise} and \texttt{echelle}, the baseline reuses the
outputs of the historical rule-based pipeline stored in
\texttt{Final\_Tagged\_Database}. The historical schema did not include a
\texttt{nature\_initiative} field. For this field, a documented deterministic
extension was applied to the first 2,500 characters of each cleaned document:
it assigns the first category whose predefined keyword rule matches the text.
This extension does not call an LLM. Its use makes the comparison possible,
but it should not be confused with a native output of the original rule-based
database.

\begin{table}[H]
\centering
\small
\begin{tabular}{lrrrrrr}
\toprule
& \multicolumn{3}{c}{Rule-based baseline} & \multicolumn{3}{c}{Local LLM} \\
\cmidrule(lr){2-4} \cmidrule(lr){5-7}
Field & Precision & Recall & F1 & Precision & Recall & F1 \\
\midrule
Target audience & 0.283 & 0.283 & 0.283 & 1.000 & 1.000 & 1.000 \\
Territorial scale & 0.347 & 0.436 & 0.386 & 1.000 & 0.987 & 0.993 \\
Initiative type & 0.589 & 0.550 & 0.569 & 1.000 & 1.000 & 1.000 \\
\bottomrule
\end{tabular}
\caption{Paired comparison between the deterministic rule-based baseline and
the local LLM on the 60-document human reference sample. Scores are micro
precision, recall, and F1.}
\label{tab:rule-based-llm-comparison}
\end{table}

The LLM achieved higher F1 scores for all three fields. The largest difference
appeared for target audience ($+0.717$ F1) and territorial scale ($+0.607$ F1).
The historical rule-based outputs for these fields were often free-text
candidate excerpts or over-inclusive keyword matches, whereas the LLM prompt
required a normalised target-audience label and a restricted list of territorial
scales. For initiative type, the keyword baseline reached an F1 of 0.569. It
correctly identified explicit terms such as \textit{atelier}, \textit{formation},
or \textit{rapport}, but it selected the first matching term and therefore
struggled when a document mentioned several possible forms of activity.

This result does not make the rule-based approach obsolete. Its decisions are
deterministic, easy to inspect, and useful for identifying explicit lexical
signals. However, on this heterogeneous corpus, the constrained LLM is better
able to select a field value that corresponds to the intended metadata schema.
The comparison is diagnostic rather than a definitive benchmark: it uses one
60-document, source-stratified sample and one human annotator. A larger blind
reference annotation and a second annotator would be needed to quantify the
generalisation of the observed differences.
